\begin{lemma}[Colors on the current tree are controlled]
For \(n=q+1\), use
\(T_n=\mathcal T.\mathrm{tree}(q+1)\),
\(\varphi_{n-1}=\mathcal T.\mathrm{coloring}(q)\), and
\(\varphi_n=\mathcal T.\mathrm{coloring}(q+1)\).  Then
\[
 \varphi_{n-1}\langle T_n\rangle
   \subseteq\overline\varphi_{n-1}(T_n)\cup D_{n-1},
 \qquad
 \varphi_n\langle T_n\rangle
   \subseteq\overline\varphi_n(T_n)\cup D_n.
\]
Moreover, if \(\sigma\) is a proper coloring leaving exactly \(e_0\)
uncolored and is \((T_n,D_n,\varphi_n)\)-stable, then
\(\sigma(g)=\varphi_n(g)\) for every displayed edge identity \(g\) of
\(T_n\), including the root identity \(e_0\).
\end{lemma}
\begin{proof}
We prove the first inclusion by induction on \(n\).  Recall from
\noderef{TreeUsedColors} that the notation
\(\psi\langle T\rangle\) records the colors of all displayed edges after
the distinguished first edge; that first edge is handled separately at
the end.

For \(n=1\), the first-edge and attachment fields in
\noderef{TreeSequence} identify the root and its displayed successors.
The initial constructor in \noderef{TashkinovSeriesStep} makes \(T_1\) a
\noderef{TashkinovClosure} of the one-edge seed \(T_0\), under
\(\varphi_0\), and gives \(D_0=\varnothing\).  Expand
\noderef{TashkinovAugmentationSequence}.  At each one-edge step,
\noderef{TashkinovAugmentation} says that the new displayed edge has a
color missing on the preceding prefix.  The prefixes are ordered initial
segments by \noderef{TreeSequenceInitialSegment}; hence their vertex sets
are nested, and \noderef{TreeMissingColors} puts that color among the
colors missing on the final tree \(T_1\).  Induction over the finite
augmentation sequence accounts for every displayed edge after the seed
edge.  Thus
\(\varphi_0\langle T_1\rangle\subseteq
\overline\varphi_0(T_1)=\overline\varphi_0(T_1)\cup D_0\).

Now let \(n>1\), and assume the assertion for \(T_{n-1}\) under
\(\varphi_{n-2}\).  Consider first an edge already displayed on
\(T_{n-1}\).  Its old color lies in
\(\overline\varphi_{n-2}(T_{n-1})\cup D_{n-2}\) by induction.  It is
incident with \(T_{n-1}\), so
\noderef{HkuIncidentColorPreserved}, applied to iteration \(n-1\), says
that its \(\varphi_{n-1}\)-color belongs to
\(\overline\varphi_{n-1}(T_{n-1})\cup D_{n-1}\).  Since
\(T_{n-1}\) is an initial segment of \(T_n\), this set is contained in
\(\overline\varphi_{n-1}(T_n)\cup D_{n-1}\).

It remains to treat the edges newly displayed between \(T_{n-1}\) and
\(T_n\).  Inspect the certificate for iteration \(n-1\) in
\noderef{TashkinovSeries}.  For a PE, \(T_n\) is obtained by a closure of
\(T_{n-1}\) under \(\varphi_{n-1}\); every new edge is therefore a
Tashkinov augmentation edge and its color is missing on an earlier prefix,
hence on \(T_n\).  For an SE or RE, the constructor first adjoins its exact
connecting edge in the sense of \noderef{TreeSequenceAdjoinEdge} and then
takes such a closure.  The connecting edge's
color is one of the exact colors stored at that iteration, and therefore
belongs to
\(\overline\varphi_{n-1}(T_n)\cup D_{n-1}\) after expanding
\noderef{TashkinovStoredColors} and using
\(A\cup(B\setminus A)=A\cup B\).  Every subsequent closure edge is again
colored by a color missing on its preceding augmentation prefix, hence on
the final tree.  These cases exhaust all newly displayed identities and
prove
\[
 \varphi_{n-1}\langle T_n\rangle
 \subseteq\overline\varphi_{n-1}(T_n)\cup D_{n-1}.               \tag{1}
\]

For every edge counted on the left of (1), its color and edge identity
satisfy the hypotheses of \noderef{HkuIncidentColorPreserved} at the
current iteration \(n\).  That lemma's unconditional conclusion gives a
\(\varphi_n\)-color in
\(\overline\varphi_n(T_n)\cup D_n\).  Thus
\[
 \varphi_n\langle T_n\rangle
 \subseteq\overline\varphi_n(T_n)\cup D_n,                       \tag{2}
\]
including the possible PE connecting-edge exception because the cited
lemma records its replacement color rather than merely its old color.

Finally let \(\sigma\) have the three properties in the statement.  If a
displayed edge \(g\) is not the first edge, the series certificate says it
is colored under \(\varphi_n\), and (2) puts that color in the controlling
set.  The edge is incident with the tree by its attachment equation in
\noderef{TreeSequence}, which is exactly \noderef{TreeIncidentEdge}.
The edge clause of \noderef{StableColoring} therefore
gives \(\sigma(g)=\varphi_n(g)\).  The first displayed edge equals
\(e_0\).  The relevant constructor of \noderef{TashkinovSeriesStep} says
that \(\varphi_n\) leaves exactly \(e_0\) uncolored, while the same is an
explicit hypothesis on \(\sigma\); by \noderef{IsOneEdgeUncolored}, both
values there are \(\mathrm{none}\).  Equality follows also for this root
edge, completing the proof.
\end{proof}
